\chapter{Daten}\label{ch:daten}

\section{Daten manipulieren}

\subsection{Python-Dictionaries}

Ein Python-Dictionary ist eine Datenstruktur, die dazu verwendet wird, um eine Sammlung von Schlüssel-Wert-Paaren zu speichern. Jeder Schlüssel muss eindeutig sein und auf einen bestimmten Wert verweisen.


\subsubsection*{Beispiel 1: Informationen zu einer Person speichern}

Ein einfaches Dictionary, das Informationen über eine Person speichert:

\begin{lstlisting}[language=Python]
person = {
    'vorname': 'Max',
    'nachname': 'Mustermann',
    'alter': 25,
    'beruf': 'Entwickler'
}

# Zugriff auf Werte über die Schlüssel
print("Vorname:", person['vorname'])
print("Alter:", person['alter'])
\end{lstlisting}

\subsubsection*{Beispiel 2: Wörter zählen in einem Satz}

Ein Dictionary, um die Häufigkeit der Wörter in einem Satz zu zählen:

\begin{lstlisting}[language=Python]
satz = "Python ist eine mächtige Programmiersprache. Python wird oft für Datenanalyse verwendet."

wort_haeufigkeit = {}
woerter = satz.split()

for wort in woerter:
    if wort in wort_haeufigkeit:
        wort_haeufigkeit[wort] += 1
    else:
        wort_haeufigkeit[wort] = 1

print("Wort Häufigkeit:", wort_haeufigkeit)
\end{lstlisting}

\subsubsection*{Beispiel 3: Ein verschachteltes Dictionary für Bücher}

Ein verschachteltes Dictionary für Informationen zu Büchern:

\begin{lstlisting}[language=Python]
buch1 = {
    'titel': 'Der Hobbit',
    'autor': 'J.R.R. Tolkien',
    'genre': 'Fantasy'
}

buch2 = {
    'titel': '1984',
    'autor': 'George Orwell',
    'genre': 'Dystopie'
}

bibliothek = {
    'buch1': buch1,
    'buch2': buch2
}

# Zugriff auf Informationen über Bücher
print("Titel von Buch 1:", bibliothek['buch1']['titel'])
print("Autor von Buch 2:", bibliothek['buch2']['autor'])
\end{lstlisting}



Übungsaufgaben
\begin{itemize}
\item \textbf{Notenbuch}: Erstelle ein Programm, das ein Notenbuch repräsentiert. Verwende ein Dictionary, um für verschiedene Schüler Noten zu speichern. Das Programm sollte in der Lage sein, Noten hinzuzufügen, zu aktualisieren, zu löschen und den Durchschnitt für jeden Schüler zu berechnen.
Übungsaufgabe: Wörterbuch übersetzen
\item \textbf{Wörterbuch}: Erstelle ein Wörterbuch mit englischen und deutschen Begriffen als Schlüssel-Wert-Paare. Schreibe eine Funktion, die einen englischen Begriff als Eingabe nimmt und die deutsche Übersetzung zurückgibt. Wenn der Begriff nicht im Wörterbuch vorhanden ist, sollte die Funktion eine entsprechende Meldung ausgeben.
\end{itemize}